% Generated by build_appendix.py: INSTRUCTIONS.md and AgentDiscussion files rendered with pandoc
\begin{mdfile}{INSTRUCTIONS.md}{Working rules and workflow, written by the authors}
\mdh{1}{We are working on the Taiwanese Bankruptcy Prediction}

\begin{itemize}
\tightlist
\item
  (\url{https://archive.ics.uci.edu/dataset/572/taiwanese+bankruptcy+prediction}), UCI Machine Learning Repository. Collected from the Taiwan Economic Journal, 1999--2009.
\end{itemize}

\mdh{2}{Our goal is to build a 4-step project that goes as follows:}

\begin{enumerate}
\def\labelenumi{\arabic{enumi}.}
\tightlist
\item
  Problem Formulation \& Exploratory Data Analysis
\item
  Variable Selection
\item
  Model Selection and Decision Tree Tuning
\item
  Evaluation of the final model
\end{enumerate}

Our end-deliverable is one notebook file (ipynb), containing the code that implements the 4 steps.

\textbf{IMPORTANT}: WE MUST CODE IN THE MOST SIMPLE WAY POSSIBLE. WHAT IS MEANT BY SIMPLE IS:

\begin{itemize}
\tightlist
\item
  A python cell only does one task at a time.
\item
  We always privilege code minimality, what can be done in two lines of codes must be done in no more than two lines, always.
\item
  We always privilege simplicity and correctness over complexity and long code that is hard to visually understand.
\item
  We always include comments in the beginning of a Python cell to explain what is being done.
\item
  All the steps we code must build on top of the previous ones.
\item
  When doing assumptions, or making decisions, they must never be done silently and appear in the form of concise markdowns cells.
\item
  Markdown cells must follow a clean and coherent structure all the way through the notebook.
\end{itemize}

\mdh{2}{In the context folder you will find two elements:}

\begin{enumerate}
\def\labelenumi{\arabic{enumi}.}
\tightlist
\item
  The lecture decks of the course that will serve as great context to scope well the work we are doing.
\item
  The official instruction file for the project, provided by our teacher. Feel free to read it to have context on the final submission but don\textquotesingle t pay too much attention to it for now as we are only working on the code in a notebook for now. Also, don\textquotesingle t pay attention to the LLM disclosure element, I will take care of it when writing the report file.
\end{enumerate}

\mdh{2}{The workflow will be as follows:}

We will be working step by step and not one-shot the full project.

\begin{enumerate}
\def\labelenumi{\arabic{enumi}.}
\tightlist
\item
  You will explain very concisely your thoughts, assumptions and rationale in a folder called AgentDiscussion, where you will add a new markdown file at the beginning of each of the 4 steps of the project.
\item
  I will review in details the markdown files and either leave them as is or bring modification to them.
\item
  You will check the last state of the markdown cell and build according to it.
\item
  After you end coding, write a recap markdown file in the same AgentDiscussion folder and wait for me to review the code and the markdown file before signaling that we can proceed with the next step of the project.
\end{enumerate}
\end{mdfile}

\begin{mdfile}{AgentDiscussion/\allowbreak{}step1/\allowbreak{}step1\_plan.md}{Step 1 plan: written by the LLM, reviewed and edited by the authors}
\mdh{1}{Step 1 --- Problem Formulation \& EDA (plan)}

\mdh{2}{1. Problem formulation}

{\def\LTcaptype{none} % do not increment counter
\begin{tabularx}{\linewidth}{@{}l>{\raggedright\arraybackslash}X@{}}
\toprule\noalign{}
Element & Proposal \\
\midrule\noalign{}
\textbf{Target} & \texttt{Bankrupt?} (1 = the firm went bankrupt under Taiwan Stock Exchange rules, 0 = it did not) \\
\textbf{Inputs} & 95 financial ratios, assumed to be known \emph{before} the outcome (so no leakage) \\
\textbf{Decision supported} & A credit analyst decides whether to \textbf{flag} a firm as high-risk (deny or reprice credit, start a closer review) or \textbf{not flag} it \\
\textbf{Actions} & a = 1 flag, a = 0 do not flag \\
\bottomrule
\end{tabularx}
}

\textbf{Loss table (Lecture 1)}. This is an assumption, not an estimate:

{\def\LTcaptype{none} % do not increment counter
\begin{tabularx}{\linewidth}{@{}l>{\raggedright\arraybackslash}X>{\raggedright\arraybackslash}X@{}}
\toprule\noalign{}
& Healthy (Y = 0) & Bankrupt (Y = 1) \\
\midrule\noalign{}
Not flag (a = 0) & 0 & \textbf{C\_FN = 10} (credit loss) \\
Flag (a = 1) & \textbf{C\_FP = 1} (unnecessary review, lost client) & 0 \\
\bottomrule
\end{tabularx}
}

\begin{itemize}
\tightlist
\item
  \textbf{Bayes decision rule:} flag a firm when P(bankrupt \textbar{} x) ≥ C\_FP / (C\_FP + C\_FN) = \textbf{1/11 ≈ 0.09}.
\item
  \textbf{Primary metric:} average misclassification cost per firm, (10·FN + 1·FP) / n. This is the empirical risk under the loss table, so lower is better.
\item
  \textbf{Secondary metrics:} the confusion matrix, recall (share of bankruptcies caught) and precision (share of flags that are real bankruptcies).
\item
  \textbf{Why not accuracy:} only 3.2\% of firms are bankrupt, so "never flag" already reaches 96.8\% accuracy.
\item
  \textbf{Reference baseline:} "never flag" (the majority class) costs about 10 × 0.032 ≈ \textbf{0.32 per firm}. Every model has to beat this.
\item
  \textbf{Sensitivity:} in Step 4 we will re-run the evaluation with cost ratios of 5 and 20, because 10 is a stated assumption.
\end{itemize}

\mdh{2}{2. Assumptions and decisions (each one becomes a short markdown cell)}

\begin{enumerate}
\def\labelenumi{\arabic{enumi}.}
\tightlist
\item
  \textbf{The data has no firm ID and no year, so firms are treated as independent draws.} A time-based split is impossible, so we use a random stratified split. This becomes a stated limitation in Step 4.
\item
  \textbf{Drop \texttt{Net\ Income\ Flag}.} It has the same value in every row, so it carries no information.
\item
  \textbf{Hold out a stratified 80/20 train/test split right after loading} (\texttt{random\_state\ =\ 42}). That gives a test set of about 1,364 firms, about 44 of them bankrupt. The test set stays locked until Step 4.
\item
  \textbf{Run the EDA on the training set only}, so that no later decision uses test information (Lecture 2: "condition on every choice made using training data").
\item
  \textbf{Leave the raw values unchanged in Step 1.} 24 columns mix values in {[}0, 1{]} with values between about 6e5 and 1e10. Some columns have only a few such values; for others it is most rows (e.g. \texttt{Total\ Asset\ Growth\ Rate}: 6,017 of 6,819). These values look like an encoding artifact. A decision tree only depends on the order of the values, so it is unaffected. A linear selector (LASSO) is affected, so its treatment is decided in Step 2 and fitted inside the cross-validation folds.
\end{enumerate}

\mdh{2}{3. Notebook structure (applies to the whole notebook)}

\begin{Verbatim}[breaklines, fontsize=\footnotesize]
# Taiwanese Bankruptcy Prediction        ← title, dataset citation, 4-step overview
## 0. Setup                              ← imports, constants (SEED, C_FN, C_FP, THRESHOLD)
## 1. Problem Formulation & EDA
### 1.1 Problem formulation
### 1.2 Data loading and integrity checks
### 1.3 Train/test split
### 1.4 Exploratory analysis (training set)
### 1.5 Step 1 takeaways
## 2. Variable Selection                 ← same sub-structure, defined in Step 2's plan
...
\end{Verbatim}

\begin{itemize}
\tightlist
\item
  Decisions and assumptions are blockquotes: \texttt{\textgreater{}\ **Decision:**\ …} or \texttt{\textgreater{}\ **Assumption:**\ …}.
\item
  Findings are one-line \texttt{**Observation:**\ …} cells placed after the output they describe.
\item
  Dataset citation: UCI ML Repository, \emph{Taiwanese Bankruptcy Prediction}, DOI 10.24432/C5004D. Original study: Liang et al. (2016), \emph{European Journal of Operational Research} 252(2).
\end{itemize}

\mdh{2}{4. Planned code cells (one task each)}

{\def\LTcaptype{none} % do not increment counter
\begin{tabularx}{\linewidth}{@{}l>{\raggedright\arraybackslash}X@{}}
\toprule\noalign{}
\# & Cell \\
\midrule\noalign{}
1 & Imports \\
2 & Constants: \texttt{SEED}, \texttt{C\_FN}, \texttt{C\_FP}, \texttt{THRESHOLD} \\
3 & Load the CSV and strip the leading spaces from column names \\
4 & Shape and first rows \\
5 & Missing values and duplicate rows \\
6 & Target counts and proportions \\
7 & Bar plot of the target \\
8 & Find the constant columns, then drop them \\
9 & Split into X and y, then stratified train/test split \\
10 & Check the bankrupt rate in train and test \\
11 & Scale check: columns with values \textgreater{} 1 (count per column) \\
12 & Redundancy: number of feature pairs with \textbar r\textbar{} \textgreater{} 0.95, plus the top pairs \\
13 & Univariate signal: ROC-AUC of each feature alone, top 10 \\
14 & Boxplots of the 4 strongest features by class \\
\bottomrule
\end{tabularx}
}

What I expect to find, from a first look: 23 feature pairs with \textbar r\textbar{} \textgreater{} 0.95, several of them exact duplicates (r = 1.00). Profitability ratios (EPS, ROA, net income / assets) separate the classes best on their own, with AUC around 0.86--0.89. Both points motivate Step 2.

\mdh{2}{5. Points for you to confirm or change}

\begin{itemize}
\tightlist
\item
  \textbf{Cost ratio C\_FN : C\_FP = 10 : 1.} Alternatives: about 30 : 1, which roughly corresponds to balancing the two classes, or a non-cost metric such as F1 or recall at a fixed precision.
\item
  \textbf{Test size of 20\%.} Moving to 25\% gives about 55 bankrupt firms in the test set instead of 44.
\item
  **EDA on the full dataset.
\end{itemize}

\mdh{2}{Environment}

The project has a \texttt{.venv} created with \texttt{uv} (Python 3.12), and the dependencies are pinned in \texttt{requirements.txt}. The notebook will be saved with its outputs, at \texttt{notebooks/bankruptcy\_prediction.ipynb}.
\end{mdfile}

\begin{mdfile}{AgentDiscussion/\allowbreak{}step1/\allowbreak{}step1\_recap.md}{Step 1 recap: written by the LLM, reviewed by the authors}
\mdh{1}{Step 1 --- Problem Formulation \& EDA (recap)}

\textbf{Notebook:} \texttt{notebooks/bankruptcy\_prediction.ipynb}, sections 0 and 1. It runs end to end on the \texttt{.venv} kernel and is saved with its outputs.

\mdh{2}{What was built}

{\def\LTcaptype{none} % do not increment counter
\begin{tabularx}{\linewidth}{@{}l>{\raggedright\arraybackslash}X@{}}
\toprule\noalign{}
Section & Content \\
\midrule\noalign{}
Title & Short introduction, dataset citation, outline of the 4 steps \\
0. Setup & Imports; constants \texttt{SEED\ =\ 42}, \texttt{C\_FN\ =\ 10}, \texttt{C\_FP\ =\ 1}, \texttt{THRESHOLD\ =\ 1/11} \\
1.1 Problem formulation & Target, inputs, decision, loss table, Bayes rule, metric, baseline, plus 2 assumption cells (costs; no leakage and independent firms) \\
1.2 Loading and integrity & Load and clean column names; shape; missing values and duplicates; class counts and bar plot; drop the constant column; separate X and y \\
1.3 Exploratory analysis & Values above 1; Spearman correlation matrix (heatmap) and correlated pairs; AUC of each feature alone; boxplots of the top 4 features \\
1.4 Train/test split & Stratified 80/20 split and a check of both sets \\
1.5 Takeaways & 5 bullets that summarise the step \\
\bottomrule
\end{tabularx}
}

There are 18 code cells and 7 blockquote cells (decisions and assumptions), plus one \texttt{**Observation:**} line after each output that needs interpreting.

\mdh{2}{Key results}

\begin{itemize}
\tightlist
\item
  \textbf{Data quality:} 6,819 firms, no missing values, no duplicate rows. \texttt{Net\ Income\ Flag} has the same value everywhere and is dropped, which leaves \textbf{94 features}.
\item
  \textbf{Class imbalance:} 220 bankrupt firms (3.2\%). The "never flag" baseline costs about \textbf{0.32 per firm}.
\item
  \textbf{Extreme values:} 24 features mix values in {[}0, 1{]} with values between 6e5 and 1e10, with nothing in between. A few firms are affected in some features; 6,017 firms are affected in \texttt{Total\ Asset\ Growth\ Rate}.
\item
  \textbf{Redundancy:} 50 pairs of features have Spearman \textbar \ensuremath{\rho}\textbar{} \textgreater{} 0.95. These pairs involve 47 of the 94 features, and 13 of the pairs have \textbar \ensuremath{\rho}\textbar{} \textgreater{} 0.999. The heatmap shows blocks of related ratios: profitability, per-share values, growth, leverage and cash flow.
\item
  \textbf{Signal:} the strongest single predictors are profitability ratios (EPS, net income / assets, ROA), with AUC between 0.86 and 0.89. Bankrupt firms clearly score lower, but the two classes still overlap.
\item
  \textbf{Split:} 5,455 firms for training (176 bankrupt) and 1,364 for testing (44 bankrupt).
\end{itemize}

\mdh{2}{Differences from the plan}

\begin{enumerate}
\def\labelenumi{\arabic{enumi}.}
\tightlist
\item
  \textbf{Exploration on the full dataset.} I followed your edit in section 5 of the plan. Two places in the plan still say "training set only": §2 item 4 and the §3 outline. I treated your section 5 edit as the decision.
\item
  \textbf{Section order is now 1.3 EDA \ensuremath{\rightarrow} 1.4 split} (the plan had split \ensuremath{\rightarrow} EDA). Since the exploration uses the full data, splitting before it would serve no purpose, and the notebook reads more naturally this way. Tell me if you want the original order back.
\item
  \textbf{The "find constant columns, then drop them" step is two cells}, with the decision cell in between. This follows the "one task per cell" rule.
\item
  \textbf{Separating X and y has its own cell at the end of 1.2}, so the exploration and the split share the same \texttt{X} and \texttt{y}.
\item
  \textbf{The boxplots hide outliers} (\texttt{showfliers=False}, and the plot title says so). With outliers shown, the boxes were unreadable.
\end{enumerate}

\mdh{2}{Changes after your review}

\begin{itemize}
\tightlist
\item
  \textbf{Added a correlation matrix} (heatmap of all 94 features, in column order) to section 1.3, just before the list of correlated pairs.
\item
  \textbf{Switched from Pearson to Spearman correlation}, for the heatmap and for the list of pairs. A decision cell in the notebook records this.

  \begin{itemize}
  \tightlist
  \item
    Why: Pearson missed ratios that are inverses of each other, such as \texttt{Current\ Ratio} and \texttt{Current\ Liability\ to\ Current\ Assets} (\ensuremath{\rho} = 0.999), and it was distorted by the extreme values.
  \item
    Effect: the count of redundant pairs went from 23 to 50.
  \item
    Spearman measures order, which is all a tree uses.
  \item
    To go back to Pearson, change \texttt{method="spearman"} in a single cell.
  \end{itemize}
\item
  \textbf{Wrote the math in LaTeX} (\texttt{\$...\$}) in every markdown cell that contains math. In 1.1, the decision rule and the cost metric are now displayed equations, and the rule includes its one-line derivation.
\item
  \textbf{Corrected the cell counts} in this recap. The first version said 21 code cells and 7 decision cells; the real numbers were 16 and 6.
\end{itemize}

\mdh{2}{Carried into Step 2}

\begin{itemize}
\tightlist
\item
  The 24 features with extreme values need treatment before any linear or LASSO selector. That treatment must be fitted inside the cross-validation folds.
\item
  Selection and tuning only see \texttt{X\_train} and \texttt{y\_train}. \texttt{X\_test} and \texttt{y\_test} are not touched until Step 4.
\item
  New imports will be added to the single Setup cell rather than scattered through the notebook.
\end{itemize}
\end{mdfile}

\begin{mdfile}{AgentDiscussion/\allowbreak{}step2/\allowbreak{}step2\_plan.md}{Step 2 plan: written by the LLM, reviewed and edited by the authors}
\mdh{1}{Step 2 --- Variable Selection (plan)}

Everything in this step uses the training set only (\texttt{X\_train}, \texttt{y\_train}).

\mdh{2}{1. Method: LASSO logistic regression}

The outcome is binary, so the LASSO of Lecture 4 becomes an \textbf{L1-penalised logistic regression}, with an unpenalised intercept:

\[\min_{b_0,\,\beta}\; -\frac{1}{n}\sum_{i=1}^{n}\Big[y_i \log p_i + (1-y_i)\log(1-p_i)\Big] + \lambda \lVert\beta\rVert_1, \qquad p_i = \frac{1}{1+e^{-(b_0 + x_i^\top\beta)}}\]

In scikit-learn the penalty is expressed as \texttt{C} \(= 1/(n\lambda)\): a smaller \texttt{C} means a stronger penalty and fewer variables.

\textbf{Why LASSO rather than forward selection (Lecture 3.2):}

\begin{itemize}
\tightlist
\item
  One parameter controls the number of variables.
\item
  It copes with the 50 highly correlated pairs found in Step 1.
\item
  One fit takes about 0.5 s, so it can be refit inside every cross-validation fold of Step 3, as the brief requires.
\item
  Forward selection is greedy and would need thousands of model fits per selection (up to 94 candidates per step), repeated again in every fold.
\end{itemize}

\textbf{Role in the pipeline:} the LASSO only \emph{chooses} the variables. The decision tree in Step 3 receives those variables with their \textbf{raw values}, so its split thresholds stay in the original ratio units.

\mdh{2}{2. Decisions (each one becomes a short markdown cell)}

\begin{enumerate}
\def\labelenumi{\arabic{enumi}.}
\item
  \textbf{The LASSO input is transformed with a rank-based quantile transform} (to a standard normal), fitted inside each fold. This transform is not passed on to the tree.

  \begin{itemize}
  \item
    Why not cap the extreme values at 1: they carry signal. Within the extreme-value group of \texttt{Total\ Asset\ Growth\ Rate}, the magnitude alone gives an AUC of 0.72. Capping would erase that.
  \item
    What the transform does: it keeps the order of the values, removes the magnitude problem, and puts all features on the same scale, which a fair L1 penalty needs.
  \item
    Evidence, 5-fold CV log-loss on the training set:

    {\def\LTcaptype{none} % do not increment counter
    \begin{tabularx}{\linewidth}{@{}l>{\raggedright\arraybackslash}X@{}}
    \toprule\noalign{}
    Transform & CV log-loss \\
    \midrule\noalign{}
    Quantile & \textbf{0.085} \\
    Cap at 1 + standardise & 0.098 \\
    log + standardise & 0.100 \\
    Standardise only & 0.100 \\
    \bottomrule
\end{tabularx}
    }
  \end{itemize}
\item
  \textbf{The intercept stays unpenalised, as in the course.}

  \begin{itemize}
  \tightlist
  \item
    The problem: scikit-learn\textquotesingle s \texttt{liblinear} solver penalises the intercept by default, which visibly distorts the path at strong penalties.
  \item
    The fix: setting \texttt{intercept\_scaling=100} makes that penalty negligible.
  \item
    Check: the \texttt{saga} solver, which never penalises the intercept, gives the same path but runs about 6× slower and raises a convergence warning.
  \end{itemize}
\item
  \textbf{No class weights in the selector.} It is a penalised maximum-likelihood fit. The misclassification costs enter in Step 3, where the decision is made.
\item
  \textbf{\texttt{C} is chosen by 5-fold stratified cross-validation, minimising log-loss.}

  \begin{itemize}
  \tightlist
  \item
    Why log-loss: it is a proper scoring rule for probabilities, the classification counterpart of the validation MSE used in Lectures 3--4.
  \item
    Grid: 13 values of \texttt{C} from \(10^{-3}\) to \(1\), evenly spaced on a log scale.
  \item
    The same CV splitter (\texttt{StratifiedKFold(5,\ shuffle=True,\ random\_state=SEED)}) is reused in Step 3.
  \end{itemize}
\item
  \textbf{Selection stability is measured with 100 stratified bootstrap resamples} of the training set (Lecture 4). \texttt{C} stays fixed and the full procedure is refit on each resample. The resulting selection frequencies are \textbf{descriptive} and do not change the selected set.
\item
  \textbf{Hand-off to Step 3:} the selection procedure is wrapped as a scikit-learn \texttt{SelectFromModel} object. It sits in front of the tree and is \textbf{refit inside every CV fold}, with \texttt{C} fixed at the value chosen here.

  \begin{itemize}
  \tightlist
  \item
    Step 3 also compares this fixed \texttt{C} with tuning \texttt{C} together with the tree, on the decision cost.
  \item
    Both versions use the same folds and the same metric.
  \item
    The joint version is preferred only if its CV cost is lower by more than one standard error across folds.
  \end{itemize}
\end{enumerate}

\mdh{2}{3. Notebook structure}

\begin{Verbatim}[breaklines, fontsize=\footnotesize]
## 2. Variable Selection
### 2.1 Method                     ← objective, why LASSO, decision cells 1–3
### 2.2 Choice of the penalty      ← CV over C, path plot
### 2.3 Selected variables         ← coefficients at the chosen C
### 2.4 Selection stability        ← bootstrap frequencies
### 2.5 Step 2 takeaways
\end{Verbatim}

\mdh{2}{4. Planned code cells (one task each)}

New imports go into the existing Setup cell.

{\def\LTcaptype{none} % do not increment counter
\begin{tabularx}{\linewidth}{@{}l>{\raggedright\arraybackslash}X@{}}
\toprule\noalign{}
\# & Cell \\
\midrule\noalign{}
1 & Define the 5-fold stratified CV splitter (reused in Step 3) \\
2 & Define the LASSO pipeline: quantile transform, then L1 logistic regression \\
3 & Cross-validate the grid of \texttt{C} values (log-loss) \\
4 & Count the selected variables for each \texttt{C} (refit on the full training set) \\
5 & Plot CV log-loss and the number of selected variables against \texttt{C}, with the chosen \texttt{C} marked \\
6 & Fit the LASSO at the chosen \texttt{C}; list the selected variables with their coefficients \\
7 & Bootstrap: selection frequency of each variable over 100 resamples \\
8 & Bar chart of the selection frequencies \\
9 & Build the selector for Step 3 (\texttt{SelectFromModel} wrapping the LASSO pipeline at the chosen \texttt{C}) \\
\bottomrule
\end{tabularx}
}

\textbf{Preview from my checks:}

\begin{itemize}
\tightlist
\item
  \textbf{Chosen penalty:} the best \texttt{C} is about 0.18. It keeps \textbf{43 of the 94 variables}, with a CV log-loss of 0.085 against about 0.14 for the intercept-only model.
\item
  \textbf{Stability:} of the selected variables, only about 18 appear in at least 80\% of resamples, and the median selection frequency is 0.65. Correlated ratios swap places from one resample to the next, which is the instability Lecture 4 warns about.
\item
  \textbf{Runtime:} Step 2 adds about 2 minutes to the notebook (40 s for the path, about 65 s for the bootstrap).
\end{itemize}

\mdh{2}{5. Points for you to confirm or change}

\begin{itemize}
\tightlist
\item
  \textbf{Rule for choosing \texttt{C}: decided.} We use the lowest CV log-loss (course protocol). Step 3 compares it with tuning \texttt{C} together with the tree (see decision 6). The one-standard-error rule was dropped because its result depends on the solver.
\item
  \textbf{Number of bootstrap resamples:} 100, which takes about 1 minute.
\end{itemize}
\end{mdfile}

\begin{mdfile}{AgentDiscussion/\allowbreak{}step2/\allowbreak{}step2\_recap.md}{Step 2 recap: written by the LLM, reviewed by the authors}
\mdh{1}{Step 2 --- Variable Selection (recap)}

\textbf{Notebook:} section 2 of \texttt{notebooks/bankruptcy\_prediction.ipynb}. The whole notebook runs end to end in about 1.5 minutes and is saved with its outputs. Step 2 uses \texttt{X\_train} and \texttt{y\_train} only.

\mdh{2}{What was built}

{\def\LTcaptype{none} % do not increment counter
\begin{tabularx}{\linewidth}{@{}l>{\raggedright\arraybackslash}X@{}}
\toprule\noalign{}
Section & Content \\
\midrule\noalign{}
Setup & New imports added to your imports cell; your \texttt{ggplot} style is kept \\
2.1 Method & The LASSO objective (LaTeX); decision LASSO vs forward selection; evidence that the extreme values carry signal; decisions on the quantile transform, the unpenalised intercept and the absence of class weights; definition of the \texttt{lasso} pipeline \\
2.2 Choice of the penalty & Decision on CV log-loss; the \texttt{cv} splitter; cross-validation over a 13-value grid of \texttt{C}; number of variables per \texttt{C}; a 2-panel plot of the path \\
2.3 Selected variables & The 43 coefficients at the chosen \texttt{C} \\
2.4 Selection stability & Decision on the bootstrap; 100 stratified resamples; bar chart of the frequencies; decision on the hand-off; the \texttt{selector} object plus a check cell \\
2.5 Takeaways & 5 bullets \\
\bottomrule
\end{tabularx}
}

Section 2 adds 11 code cells and 7 decision cells.

\mdh{2}{Key results}

\begin{itemize}
\tightlist
\item
  \textbf{Extreme values carry signal.} Among the 4,816 training firms whose \texttt{Total\ Asset\ Growth\ Rate} exceeds 1, the value alone gives an AUC of 0.72 (0.28 with the direction flipped). This confirms that capping them at 1 would lose information.
\item
  \textbf{Chosen penalty:} \textbf{C ≈ 0.178}, which keeps \textbf{43 of the 94 variables}. Its CV log-loss is 0.0855, against 0.143 with no variable. The curve is flat: every value from 21 to 59 variables is within one standard error (0.002) of the best.
\item
  \textbf{What was selected:} leverage, profitability, liquidity and activity ratios. For the leading variables the signs make sense: more debt raises the risk, and more profit or liquidity lowers it.
\item
  \textbf{Stability is moderate.}

  \begin{itemize}
  \tightlist
  \item
    Only 15 of the 43 selected variables appear in at least 80\% of resamples. The median frequency is 0.69 and the lowest is 0.25.
  \item
    \texttt{Research\ and\ development\ expense\ rate} is not selected on the full training set, yet appears in 81\% of resamples.
  \item
    The LASSO keeps both members of a practically identical pair (\texttt{Current\ Liabilities/Liability} and \texttt{Current\ Liability\ to\ Liability}), each with a tiny coefficient.
  \item
    The variable list is a predictive shortlist, not a set of causes.
  \end{itemize}
\item
  \textbf{Check passed:} \texttt{selector}, fitted on the full training set, keeps exactly the 43 variables listed in 2.3.
\end{itemize}

\mdh{2}{Differences from the plan}

\begin{enumerate}
\def\labelenumi{\arabic{enumi}.}
\tightlist
\item
  \textbf{Extra evidence cell in 2.1} (the AUC within the extreme-value group). It backs up the quantile-transform decision inside the notebook. The comparison of the four transforms (0.085 vs 0.098--0.100) is \emph{cited} from my preliminary check rather than re-run, because re-running it would add about 6 minutes to the notebook.
\item
  \textbf{"LASSO rather than forward selection" is its own decision cell}, instead of a paragraph inside the method text.
\item
  \textbf{Extra check cell after the selector}, to confirm it reproduces the 43 variables.
\item
  \textbf{The cross-validation over \texttt{C} runs in parallel} (\texttt{n\_jobs=-1}). This brings the whole notebook down to about 1.5 minutes.
\item
  \textbf{The stability chart shows all 92 variables that were selected at least once.} It is tall. If you want a more compact figure for the report, I can limit it to the top 40.
\end{enumerate}

\mdh{2}{Changes after your review}

\begin{itemize}
\tightlist
\item
  \textbf{The stability chart has a proper legend.} Blue means "in the 43-variable shortlist" and grey means "not in the shortlist". A dashed line marks 80\% of resamples, the boundary for robust selection.
\item
  \textbf{The chart title is now "Bootstrap selection frequency (100 resamples)".} This needed \texttt{from\ matplotlib.patches\ import\ Patch} in the Setup cell.
\item
  \textbf{The observation cell points to the dashed line.}
\item
  \textbf{Your explanation cell under the chart is unchanged.}
\end{itemize}

\mdh{2}{Carried into Step 3}

\begin{itemize}
\tightlist
\item
  \texttt{cv} (the 5 stratified folds) and \texttt{selector} (unfitted; it is refit inside every fold) are ready.
\item
  The comparison of fixed \texttt{C} against \texttt{C} tuned together with the tree, using the decision rule agreed on: prefer the joint version only if its CV cost is lower by more than one standard error.
\item
  The misclassification costs have to enter the tree, through class weights or the threshold of 1/11. That is a decision for the Step 3 plan.
\end{itemize}
\end{mdfile}

\begin{mdfile}{AgentDiscussion/\allowbreak{}step3/\allowbreak{}step3\_plan.md}{Step 3 plan: written by the LLM, reviewed and edited by the authors}
\mdh{1}{Step 3 --- Model Selection and Decision Tree Tuning (plan)}

Everything in this step uses the training set only, with the 5 folds \texttt{cv} from Step 2. The test set is first used in Step 4.

\mdh{2}{1. Decision rule and metric}

\begin{itemize}
\tightlist
\item
  \textbf{The tree estimates probabilities and the loss table decides (Lectures 1 and 3).}

  \begin{itemize}
  \tightlist
  \item
    A standard Gini tree is grown.
  \item
    Each leaf estimates \(\hat\eta(x) = P(Y=1 \mid x)\) by its share of bankrupt firms.
  \item
    A firm is flagged when \(\hat\eta(x) \geq 1/11\), the Bayes threshold from 1.1.
  \item
    In code, scikit-learn\textquotesingle s \texttt{FixedThresholdClassifier(...,\ threshold=THRESHOLD)} wraps the model, so \texttt{predict()} returns the cost-optimal decision directly.
  \end{itemize}
\item
  \textbf{The metric is the average cost per firm}, \((C_{FN}\cdot\text{FN} + C_{FP}\cdot\text{FP})/n\). It is turned into a scikit-learn scorer and used for all tuning and comparisons.
\end{itemize}

\mdh{2}{2. Decisions (each one becomes a short markdown cell)}

\begin{enumerate}
\def\labelenumi{\arabic{enumi}.}
\item
  \textbf{The costs enter through the threshold, not through class weights.}

  \begin{itemize}
  \tightlist
  \item
    The alternative: grow the tree with cost-weighted Gini (\texttt{class\_weight=\{0:\ 1,\ 1:\ 10\}}).
  \item
    What I found: that gives the same CV cost within one standard error (0.162 vs 0.159 in my check).
  \item
    Why the threshold: it keeps the course\textquotesingle s reading of leaves as probability estimates. It also lets Step 4 re-run the cost-ratio sensitivity by changing only the threshold, with no retraining.
  \end{itemize}
\item
  \textbf{Tuning grid (Gini criterion, \texttt{random\_state\ =\ SEED}):}

  {\def\LTcaptype{none} % do not increment counter
  \begin{tabularx}{\linewidth}{@{}l>{\raggedright\arraybackslash}X@{}}
  \toprule\noalign{}
  Parameter & Values \\
  \midrule\noalign{}
  \texttt{max\_depth} & 2, 3, 4, 5, 6, 8, None (no limit) \\
  \texttt{min\_samples\_leaf} & 1, 10, 20, 40, 80, 120, 160, 200 \\
  \bottomrule
\end{tabularx}
  }

  Why the leaf sizes go so high:

  \begin{itemize}
  \tightlist
  \item
    Leaf shares are only trusted as probabilities if leaves are large enough (Lecture 3: "small leaves make probability estimates unstable").
  \item
    With a 3.2\% bankruptcy rate, a leaf of 120 firms contains only about 4 bankruptcies on average.
  \item
    My first grid stopped at 80, and that edge value won, so I extended the grid.
  \end{itemize}
\item
  \textbf{Four candidates, all cross-validated on the same \texttt{cv} folds.} The selection step is refit inside every fold.

  {\def\LTcaptype{none} % do not increment counter
  \begin{tabularx}{\linewidth}{@{}l>{\raggedright\arraybackslash}X>{\raggedright\arraybackslash}X@{}}
  \toprule\noalign{}
  & Candidate & Tuned \\
  \midrule\noalign{}
  M0 & Baseline: never flag & -- \\
  M1 & Tree on all 94 features & depth, leaf \\
  M2 & LASSO selector (\texttt{C} fixed at 0.178, from Step 2) + tree & depth, leaf \\
  M3 & LASSO selector + tree, with \texttt{C} tuned jointly & \texttt{C}, depth, leaf \\
  \bottomrule
\end{tabularx}
  }

  For M3, \texttt{C} takes the 7 values of Step 2\textquotesingle s grid from 0.01 to 0.316. That range includes 0.178, so M3 contains M2 as one of its options.
\item
  \textbf{Selection rule (as agreed in Step 2):} M2, the course protocol, is the reference. An alternative replaces it only if its CV cost is lower by more than one standard error. Standard error = standard deviation across the 5 folds / \ensuremath{\surd}5.
\item
  \textbf{The final model} is the chosen candidate refit on the full training set with its best parameters. It is then \textbf{frozen} for Step 4.
\end{enumerate}

\mdh{2}{3. Notebook structure}

\begin{Verbatim}[breaklines, fontsize=\footnotesize]
## 3. Model Selection and Decision Tree Tuning
### 3.1 Decision rule and metric        ← cost function, scorer
### 3.2 Candidates and tuning grid      ← decision cells 1–4, grid
### 3.3 Cross-validated tuning          ← M0, M1, M2, M3, validation curve
### 3.4 Model comparison                ← table, plot, selection rule
### 3.5 Final model                     ← frozen pipeline, what it uses
### 3.6 Step 3 takeaways
\end{Verbatim}

\mdh{2}{4. Planned code cells (one task each)}

New imports go into the Setup cell: \texttt{DecisionTreeClassifier}, \texttt{FixedThresholdClassifier}, \texttt{DummyClassifier}, \texttt{Pipeline}, \texttt{make\_scorer}, \texttt{confusion\_matrix}, \texttt{cross\_val\_score}.

{\def\LTcaptype{none} % do not increment counter
\begin{tabularx}{\linewidth}{@{}l>{\raggedright\arraybackslash}X@{}}
\toprule\noalign{}
\# & Cell \\
\midrule\noalign{}
1 & Cost function: average misclassification cost per firm \\
2 & Scorer built from the cost function \\
3 & Tree grid (\texttt{max\_depth}, \texttt{min\_samples\_leaf}) \\
4 & M0: CV cost of "never flag" \\
5 & M1: grid search, thresholded tree on all features \\
6 & M2: grid search, thresholded pipeline (selector with fixed \texttt{C}, then tree) \\
7 & M3: same as M2, with \texttt{C} added to the grid \\
8 & Validation curve for M2: training and CV cost against depth, for leaf size 1 and for the chosen leaf size \\
9 & Comparison table: CV cost ± standard error, number of variables, best parameters (M0--M3) \\
10 & Bar chart of the comparison with ±1 standard-error bars \\
11 & Apply the selection rule and keep the final model \\
12 & Describe the final model: parameters, variables selected, variables the tree actually uses, number of leaves \\
\bottomrule
\end{tabularx}
}

\mdh{2}{5. Preview from my checks (same folds and seed, so the notebook should reproduce them)}

{\def\LTcaptype{none} % do not increment counter
\begin{tabularx}{\linewidth}{@{}l>{\raggedright\arraybackslash}X>{\raggedright\arraybackslash}X@{}}
\toprule\noalign{}
Candidate & CV cost per firm & Best parameters \\
\midrule\noalign{}
M0 never flag & 0.323 ± 0.002 & -- \\
M1 all features & 0.163 ± 0.009 & depth 4, leaf 120 \\
\textbf{M2 LASSO, fixed \texttt{C}} & \textbf{0.151 ± 0.008} & depth 4, leaf 120 (43 selected; the tree uses 8) \\
M3 LASSO, joint \texttt{C} & 0.148 ± 0.009 & \texttt{C} = 0.01 (about 7 variables), depth 3, leaf 10 \\
\bottomrule
\end{tabularx}
}

\begin{itemize}
\tightlist
\item
  \textbf{All trees roughly halve the baseline cost.}
\item
  \textbf{Selection helps:} M2 beats M1 by 0.012, which is more than one standard error.
\item
  \textbf{M3 does not beat M2 by enough:} its gain of 0.003 is well below one standard error, so the rule keeps \textbf{M2}.

  \begin{itemize}
  \tightlist
  \item
    Interesting side result: tuning \texttt{C} on the decision cost picks a much sparser selection (about 7 variables) at practically the same cost.
  \item
    In practice the final M2 tree uses only 8 variables anyway, so the two models are similarly interpretable.
  \end{itemize}
\item
  \textbf{The validation curve tells the overfitting story (Lecture 2):} with leaves of size 1, training cost falls to 0 as depth grows, while CV cost rises from 0.17 to 0.23. With leaves of 120, depth 4 is best.
\item
  \textbf{Runtime:} Step 3 adds about 2.5 minutes, almost all of it for M3.
\end{itemize}

\mdh{2}{6. Points for you to confirm or change}

\textbf{Decided (you approved the recommendations):}

\begin{itemize}
\tightlist
\item
  The costs enter through the threshold.
\item
  The grid stays as listed in decision 2.
\item
  The agreed rule stays as it is, with no tie-breaker for fewer variables.
\end{itemize}

The original points are kept below for reference.

\begin{itemize}
\tightlist
\item
  \textbf{How costs enter:} through the threshold (recommended), or through class weights.
\item
  \textbf{The grid in decision 2.}
\item
  \textbf{Keeping the agreed rule, even though M3 reaches the same cost with about 7 variables.} The alternative is to add a tie-breaker: "within one standard error, prefer fewer selected variables", which would pick M3. I recommend keeping the agreed rule. It was set before seeing these results, and changing a rule after seeing the results weakens it.
\end{itemize}
\end{mdfile}

\begin{mdfile}{AgentDiscussion/\allowbreak{}step3/\allowbreak{}step3\_recap.md}{Step 3 recap: written by the LLM, reviewed by the authors}
\mdh{1}{Step 3 --- Model Selection and Decision Tree Tuning (recap)}

\textbf{Notebook:} section 3 of \texttt{notebooks/bankruptcy\_prediction.ipynb}. The whole notebook runs end to end in about 4 minutes; tuning M3 takes about 2 of them. It is saved with its outputs, and Step 3 uses the training set only.

\mdh{2}{What was built}

{\def\LTcaptype{none} % do not increment counter
\begin{tabularx}{\linewidth}{@{}l>{\raggedright\arraybackslash}X@{}}
\toprule\noalign{}
Section & Content \\
\midrule\noalign{}
Setup & New imports added to the single imports cell \\
3.1 Decision rule and metric & Leaf share as \(\hat\eta(x)\), flag when ≥ 1/11, empirical-risk formula (LaTeX); decision on threshold vs class weights; \texttt{avg\_cost} and \texttt{cost\_scorer} \\
3.2 Candidates and grid & Decisions on the grid, the 4 candidates and the selection rule; \texttt{tree\_grid} \\
3.3 Cross-validated tuning & M0, M1, M2, M3, each with its own cell; M2 validation curve (2 panels) \\
3.4 Model comparison & Comparison table, bar chart with ±1 standard-error bars, rule table, decision cell \\
3.5 Final model & \texttt{final\_model} frozen; description (variables selected and used, depth, leaves) \\
3.6 Takeaways & 6 bullets \\
\bottomrule
\end{tabularx}
}

Section 3 adds 13 code cells and 5 decision cells.

\mdh{2}{Key results (identical to the plan\textquotesingle s preview)}

{\def\LTcaptype{none} % do not increment counter
\begin{tabularx}{\linewidth}{@{}l>{\raggedright\arraybackslash}X>{\raggedright\arraybackslash}X>{\raggedright\arraybackslash}X>{\raggedright\arraybackslash}X>{\raggedright\arraybackslash}X>{\raggedright\arraybackslash}X@{}}
\toprule\noalign{}
Candidate & CV cost per firm & SE & Depth & Leaf & \texttt{C} & Variables \\
\midrule\noalign{}
M0 never flag & 0.3226 & 0.0016 & -- & -- & -- & 0 \\
M1 all features & 0.1630 & 0.0090 & 4 & 120 & -- & 94 \\
\textbf{M2 LASSO, \texttt{C} fixed} & \textbf{0.1512} & 0.0077 & 4 & 120 & 0.178 & 43 \\
M3 LASSO, \texttt{C} joint & 0.1483 & 0.0087 & 3 & 10 & 0.01 & 7 \\
\bottomrule
\end{tabularx}
}

\begin{itemize}
\tightlist
\item
  \textbf{M2 is kept by the agreed rule.}

  \begin{itemize}
  \tightlist
  \item
    M3\textquotesingle s gain is 0.003, well below its standard error of 0.009.
  \item
    M1 is worse by 0.012, more than one standard error, so the variable selection helps.
  \end{itemize}
\item
  \textbf{Final model:} the LASSO keeps 43 variables, followed by a depth-4 tree with \textbf{9 leaves that uses only 8 variables}. Those 8 are profitability (EPS, continuous interest rate), debt and its cost (borrowing dependency, interest-bearing debt rate, interest coverage), growth (net value growth) and activity (inventory/working capital, quick asset turnover).
\item
  \textbf{Validation curve:}

  \begin{itemize}
  \tightlist
  \item
    With leaves of 1 firm, training cost falls to 0 as depth grows, while CV cost rises from 0.17 to 0.23. That is clear overfitting.
  \item
    With leaves of at least 120 firms, the gap between training and CV cost stays small, and depth 4 is best.
  \end{itemize}
\end{itemize}

\mdh{2}{Differences from the plan}

\begin{enumerate}
\def\labelenumi{\arabic{enumi}.}
\tightlist
\item
  \textbf{M2\textquotesingle s grid states \texttt{C} explicitly} (\texttt{{[}best\_C{]}}) instead of relying on the value preset inside \texttt{selector}. This is more robust and makes M2 and M3 directly parallel, since M3 just swaps in \texttt{C\_grid{[}4:11{]}}.
\item
  \textbf{"Apply the rule and keep the final model" became two cells} (the rule table, then the freeze), to keep one task per cell. That is why there are 13 code cells instead of 12.
\item
  \textbf{Section 3.1 shows the metric as an empirical-risk formula}, and \texttt{avg\_cost} computes exactly that: the mean of the per-firm losses.
\item
  \textbf{Two pandas 3 adjustments:}

  \begin{itemize}
  \tightlist
  \item
    In the validation-curve plot, depth "None" is converted to the label "no limit" before plotting.
  \item
    The comparison table uses \texttt{convert\_dtypes()}, so depths and leaf sizes display as whole numbers.
  \end{itemize}
\end{enumerate}

\mdh{2}{Carried into Step 4}

\begin{itemize}
\tightlist
\item
  \textbf{What Step 4 uses:}

  \begin{itemize}
  \tightlist
  \item
    \texttt{final\_model}, frozen: a \texttt{FixedThresholdClassifier} whose \texttt{predict()} returns the flag decision.
  \item
    \texttt{avg\_cost}.
  \item
    The never-flag baseline, for comparison on the test set.
  \end{itemize}
\item
  \textbf{First use of the test set:}

  \begin{itemize}
  \tightlist
  \item
    test cost compared with the baseline
  \item
    confusion matrix, recall and precision
  \item
    drawing the tree for interpretation
  \end{itemize}
\item
  \textbf{Cost-ratio sensitivity (5 and 20):} \texttt{final\_model.set\_params(threshold=...)} changes the threshold (1/6, 1/21) \textbf{without refitting}.
\item
  \textbf{Limitations to discuss:}

  \begin{itemize}
  \tightlist
  \item
    no time-based split
  \item
    the costs are stipulated
  \item
    CV noise (SE about 0.008, with only about 35 bankruptcies per fold)
  \item
    no nested CV for the M2 vs M3 comparison
  \item
    moderate stability of the LASSO selection
  \end{itemize}
\end{itemize}
\end{mdfile}

\begin{mdfile}{AgentDiscussion/\allowbreak{}step4/\allowbreak{}step4\_plan.md}{Step 4 plan: written by the LLM, reviewed and edited by the authors}
\mdh{1}{Step 4 --- Evaluation of the Final Model (plan)}

This step is the \textbf{first and only use of the test set} (1,364 firms, 44 bankrupt). Following Lecture 2, the frozen final model from Step 3 is evaluated once, and nothing is tuned or re-chosen afterwards.

\textbf{I wrote this plan without computing any test-set result.} That way, the choice of what to report cannot be influenced by the results. The previews below come from the training set only.

\mdh{2}{1. Decisions (each one becomes a short markdown cell)}

\begin{enumerate}
\def\labelenumi{\arabic{enumi}.}
\tightlist
\item
  \textbf{Only the final model (M2) and the "never flag" baseline are scored on the test set.} M1 and M3 are not.

  \begin{itemize}
  \tightlist
  \item
    Reporting the test scores of rejected candidates invites choosing again on the test set, which would make the test score part of model selection (Lecture 2).
  \end{itemize}
\item
  \textbf{The test results are reported as:}

  \begin{itemize}
  \tightlist
  \item
    the average cost per firm (primary metric)
  \item
    recall, precision and the number of flagged firms
  \item
    the confusion matrix
  \item
    accuracy, shown only to illustrate why it was rejected in 1.1
  \end{itemize}
\item
  \textbf{Uncertainty: a bootstrap confidence interval, because Lecture 2\textquotesingle s Hoeffding bound is too loose here.}

  \begin{itemize}
  \tightlist
  \item
    The loss lies in {[}0, 10{]} and there are 1,364 test firms. Hoeffding then gives a margin of \textbf{±0.37} at 95\%. That is larger than the gap we want to measure (≈ 0.32 vs ≈ 0.15), so the bound is valid but uninformative (Lecture 2: "bounds larger than the loss range are valid but uninformative"). It is shown in one line.
  \item
    Instead: a \textbf{95\% percentile bootstrap} with 1,000 resamples of the test firms, keeping the predictions fixed. The interval is computed for the model\textquotesingle s cost, the baseline\textquotesingle s cost and their \textbf{paired difference}.
  \end{itemize}
\item
  \textbf{Sensitivity to the cost ratio (announced in 1.1):} \(C_{FN}/C_{FP} \in \{5, 10, 20\}\).

  \begin{itemize}
  \tightlist
  \item
    For each ratio \(r\), the decision threshold becomes \(1/(1+r)\) (Bayes rule). The tree is \textbf{not retrained}.
  \item
    Each version is compared with the baseline under the same ratio.
  \item
    In code, the probabilities are computed once (\texttt{predict\_proba}) and thresholded per ratio, so \texttt{final\_model} is never modified.
  \end{itemize}
\item
  \textbf{Switching points:} for each leaf, the cost ratio at which its decision flips, \(r^* = (1-\hat\eta)/\hat\eta\), computed from its training share \(\hat\eta\). This explains the sensitivity results directly from the loss table (Lecture 1, "different costs can reverse the decision").
\item
  \textbf{Interpretation:}

  \begin{itemize}
  \tightlist
  \item
    \textbf{Tree diagram.} The final tree is drawn with \texttt{plot\_tree}, in raw ratio units, with node ids.
  \item
    \textbf{Leaf table.} Each leaf gets its training firms, its bankruptcy share and its decision (flag if the share is ≥ 1/11).
  \item
    \textbf{Why the table is needed.} \texttt{plot\_tree} labels each leaf with its \emph{majority} class, which is "healthy" for all 9 leaves. That label would mislead, so the decision comes from the table instead.
  \end{itemize}
\item
  \textbf{The notebook closes with a decision memo} following Lecture 1\textquotesingle s structure: the action and the information used, the probability and loss assumptions, the comparison with a feasible alternative, and the limitations and conditions that would change the choice.
\end{enumerate}

\mdh{2}{2. One small change to Step 3 code}

\texttt{avg\_cost} gets optional cost arguments: \texttt{def\ avg\_cost(y\_true,\ y\_pred,\ c\_fn=C\_FN,\ c\_fp=C\_FP)}.

\begin{itemize}
\tightlist
\item
  Step 3 is unaffected, because the defaults are the same values.
\item
  The sensitivity analysis can then reuse the function instead of repeating the formula.
\end{itemize}

\mdh{2}{3. Notebook structure}

\begin{Verbatim}[breaklines, fontsize=\footnotesize]
## 4. Evaluation of the Final Model
### 4.1 Test-set performance           ← decisions 1–2, metrics table, confusion matrix
### 4.2 Uncertainty of the test estimate ← Hoeffding (1 line), bootstrap intervals
### 4.3 Sensitivity to the cost ratio  ← ratios 5/10/20, switching points
### 4.4 Interpretation of the final tree ← tree diagram, leaf table
### 4.5 Recommendation and limitations ← decision memo (markdown)
\end{Verbatim}

\mdh{2}{4. Planned code cells (one task each)}

New imports go into the Setup cell: \texttt{recall\_score}, \texttt{precision\_score}, \texttt{accuracy\_score}, \texttt{ConfusionMatrixDisplay}, \texttt{plot\_tree}.

{\def\LTcaptype{none} % do not increment counter
\begin{tabularx}{\linewidth}{@{}l>{\raggedright\arraybackslash}X@{}}
\toprule\noalign{}
\# & Cell \\
\midrule\noalign{}
1 & Helper \texttt{test\_metrics(y\_true,\ y\_pred,\ c\_fn)}: cost, recall, precision, accuracy, flagged firms \\
2 & Test predictions of the final model and of the baseline (flags nobody) \\
3 & Metrics table: final model vs baseline \\
4 & Confusion matrix of the final model (plot) \\
5 & Hoeffding half-width for the test cost (one line) \\
6 & Bootstrap: 1,000 resamples of the test firms; costs of model, baseline and difference \\
7 & 95\% percentile intervals table \\
8 & Test probabilities of the final model (computed once) \\
9 & Sensitivity table for ratios 5, 10, 20: threshold, model cost, baseline cost, recall, precision \\
10 & Leaf table: training firms, bankruptcy share, decision at 1/11, switching ratio \(r^*\) \\
11 & Plot of the final tree (node ids, proportions, raw units) \\
\bottomrule
\end{tabularx}
}

\mdh{2}{5. Preview from the training set only}

{\def\LTcaptype{none} % do not increment counter
\begin{tabularx}{\linewidth}{@{}l>{\raggedright\arraybackslash}X>{\raggedright\arraybackslash}X>{\raggedright\arraybackslash}X>{\raggedright\arraybackslash}X>{\raggedright\arraybackslash}X@{}}
\toprule\noalign{}
Leaf share \(\hat\eta\) & Training firms & Flagged at 1/6 (r = 5) & 1/11 (r = 10) & 1/21 (r = 20) & \(r^*\) \\
\midrule\noalign{}
0.420 & 150 & \ensuremath{\checkmark} & \ensuremath{\checkmark} & \ensuremath{\checkmark} & 1.4 \\
0.397 & 121 & \ensuremath{\checkmark} & \ensuremath{\checkmark} & \ensuremath{\checkmark} & 1.5 \\
0.116 & 146 & -- & \ensuremath{\checkmark} & \ensuremath{\checkmark} & 7.6 \\
0.093 & 129 & -- & \ensuremath{\checkmark} & \ensuremath{\checkmark} & 9.8 \\
0.093 & 129 & -- & \ensuremath{\checkmark} & \ensuremath{\checkmark} & 9.8 \\
0.039 & 155 & -- & -- & -- & 24.8 \\
0.016, 0.015, 0.003 & 122, 264, 4,239 & -- & -- & -- & \textgreater{} 60 \\
\bottomrule
\end{tabularx}
}

\begin{itemize}
\tightlist
\item
  \textbf{Ratios 10 and 20 give identical decisions.} Every ratio between about 9.8 and 24.8 does, so the recommendation is robust to doubling the cost of a missed bankruptcy.
\item
  \textbf{Ratio 5 flags only the two high-risk leaves:} 271 training firms instead of 675.
\item
  \textbf{Two leaves (258 firms) are flagged only narrowly.} Their share of 0.093 is barely above 1/11, so with a cost ratio below 9.8 they would not be flagged. This goes into the limitations.
\item
  \textbf{Runtime:} Step 4 adds only a few seconds.
\end{itemize}

\mdh{2}{6. Points for you to confirm or change}

\begin{itemize}
\tightlist
\item
  \textbf{Score only M2 and the baseline on the test set: decided.}
\item
  \textbf{Bootstrap: decided.} 1,000 resamples of the test firms, not stratified, so that the share of bankrupt firms also varies.
\item
  \textbf{Switching-point column (\(r^*\)) in the leaf table: decided.} Kept, as an addition to the original Step 1 plan, which only mentioned ratios 5 and 20.
\end{itemize}
\end{mdfile}

\begin{mdfile}{AgentDiscussion/\allowbreak{}step4/\allowbreak{}step4\_recap.md}{Step 4 recap: written by the LLM, reviewed by the authors}
\mdh{1}{Step 4 --- Evaluation of the Final Model (recap)}

\textbf{Notebook:} section 4 of \texttt{src/bankruptcy\_prediction.ipynb}, which you moved from \texttt{notebooks/}. I applied the edits there and kept your \texttt{.venv} kernel selection. The full notebook runs end to end in about 4 minutes with no errors and is saved with its outputs.

\textbf{The test set was used exactly once}, on the frozen final model (M2) and the "never flag" baseline only.

\mdh{2}{What was built}

{\def\LTcaptype{none} % do not increment counter
\begin{tabularx}{\linewidth}{@{}l>{\raggedright\arraybackslash}X@{}}
\toprule\noalign{}
Section & Content \\
\midrule\noalign{}
Setup & New imports: \texttt{recall\_score}, \texttt{precision\_score}, \texttt{accuracy\_score}, \texttt{ConfusionMatrixDisplay}, \texttt{plot\_tree} \\
Step 3 & \texttt{avg\_cost} gained optional \texttt{c\_fn}/\texttt{c\_fp} arguments, with the same defaults, so Step 3 results are unchanged (M2 still 0.1512) \\
4.1 Test-set performance & 2 decision cells; \texttt{test\_metrics} helper; predictions; metrics table; confusion matrix \\
4.2 Uncertainty & Decision; Hoeffding half-width; 1,000-resample bootstrap; 95\% intervals \\
4.3 Sensitivity & Decision; test probabilities computed once; table for ratios 5/10/20; decision on switching points; leaf table with \(r^*\) \\
4.4 Interpretation & Decision (why no class labels); tree diagram with node ids \\
4.5 Recommendation and limitations & Decision memo in Lecture 1\textquotesingle s structure; 7 limitations \\
\bottomrule
\end{tabularx}
}

Section 4 adds 11 code cells and 6 decision cells. The notebook now has 132 cells in total.

\mdh{2}{Key results}

{\def\LTcaptype{none} % do not increment counter
\begin{tabularx}{\linewidth}{@{}l>{\raggedright\arraybackslash}X>{\raggedright\arraybackslash}X@{}}
\toprule\noalign{}
Test set (1,364 firms, 44 bankrupt) & Final model & Baseline \\
\midrule\noalign{}
Average cost per firm & \textbf{0.179} & 0.323 \\
95\% bootstrap interval & {[}0.134, 0.233{]} & {[}0.235, 0.418{]} \\
Recall & 72.7\% (32 of 44) & 0\% \\
Precision & 20.5\% & -- \\
Flagged firms & 156 (about 1 in 9) & 0 \\
Accuracy & 90.0\% & 96.8\% \\
\bottomrule
\end{tabularx}
}

\begin{itemize}
\item
  \textbf{The model cuts the cost by 45\%.} The paired interval for the gain, {[}\ensuremath{-}0.228, \ensuremath{-}0.065{]}, excludes 0.
\item
  \textbf{The test cost (0.179) is higher than the CV cost (0.151),} but the CV value lies inside the test interval. The gap is consistent with the noise you\textquotesingle d expect from only 44 test bankruptcies.
\item
  \textbf{Hoeffding\textquotesingle s margin is ±0.37:} valid, but too wide to separate the two rules. This is what motivated using the bootstrap.
\item
  \textbf{Sensitivity:} the model beats the baseline at every cost ratio, cutting the cost by 35\% at r = 5, 45\% at r = 10 and 59\% at r = 20. Ratios 10 and 20 give identical decisions, and ratio 5 flags 68 firms.
\item
  \textbf{Switching points:}

  \begin{itemize}
  \tightlist
  \item
    Nodes 1 and 13 are flagged for any ratio above 1.5.
  \item
    Node 12 switches at 7.6.
  \item
    Nodes 6 and 15 switch at 9.75, so they are flagged only narrowly.
  \item
    No leaf switches between 9.75 and 24.8.
  \end{itemize}
\item
  \textbf{Interpretation.} The tree raises an alarm in three situations:

  \begin{itemize}
  \tightlist
  \item
    the lowest net value growth (42\% bankrupt);
  \item
    high borrowing combined with weak profitability, especially when the debt is expensive (40\% and 12\%);
  \item
    two smaller pockets at about 9\%.
  \end{itemize}

  Node 9 holds 78\% of the firms and has a bankruptcy share of only 0.3\%.
\end{itemize}

\mdh{2}{Differences from the plan}

\begin{enumerate}
\def\labelenumi{\arabic{enumi}.}
\tightlist
\item
  \textbf{The leaf table sits at the end of 4.3,} where it explains the sensitivity results, rather than in 4.4. Section 4.4 refers back to it to read the tree diagram.
\item
  \textbf{The tree diagram uses \texttt{precision=5}.} With 3 decimals, splits such as \texttt{Net\ Value\ Growth\ Rate\ \textless{}=\ 0.00037} displayed as \texttt{\textless{}=\ 0.0}. Colours and class labels are left out, as planned.
\item
  \textbf{Both results tables have one row per model or per ratio,} so counts display as whole numbers (156, not 156.0000).
\item
  \textbf{The notebook\textquotesingle s location changed} (\texttt{notebooks/} \ensuremath{\rightarrow} \texttt{src/}) partway through this step. The first run had already been written before the move, so I re-applied the final edits to \texttt{src/}.
\end{enumerate}

\mdh{2}{Notes}

\begin{itemize}
\tightlist
\item
  \texttt{step4\_plan.md} is still at the root of \texttt{AgentDiscussion/}. I put this recap in \texttt{AgentDiscussion/step4/} to match your new layout, and I left the plan where it is in case you have it open. Move it whenever you like.
\item
  The recap files for Steps 1--3 still mention \texttt{notebooks/} as the notebook\textquotesingle s location. I left them unchanged, since they record what was true at the time.
\end{itemize}
\end{mdfile}
